%************************************************
\chapter{Selezione semplice}
%************************************************

Prima di arrivare ai problemi con la selezione trovo sempre utile fare qualche esercizio sulle condizioni logiche che sono necessarie, anzi fondamentali, in tutti gli esercizi da qui in avanti.

\section{Condizioni logiche}

\begin{exercise}
	Scrivere le condizioni logiche per le seguenti situazioni:
	\begin{enumerate}
		\item la condizione è vera quando la variabile \texttt{x} è maggiore o uguale di 16;
		\item la condizione è vera quando la variabile \texttt{y} è maggiore di 50;
		\item la condizione è falsa quando la variabile \texttt{z} è maggiore di 4;
		\item la condizione è falsa quando la variabile \texttt{a} è minore di 10;
		\item la condizione è vera quando la variabile \texttt{b} è positiva.
	\end{enumerate}
\end{exercise}

\section{Selezione semplice}

Bla bla bla...

\section{Selezione meno semplice}

Bla bla bla...